% 2015-TS-ISN-ApprendreLaTeX.tex
% Mise à jour : 02/04/2016
% 
\documentclass[a4paper,11pt,french]{article}
\usepackage[utf8]{inputenc}    
\usepackage[T1]{fontenc}
\usepackage[francais]{babel}
\usepackage{graphicx}
\usepackage[top=3cm, bottom=3cm, left=3cm, right=3cm]{geometry}
\title{Compte Rendu : Bomberman}
\author{Tanguy \bsc{Thomas}, Titouan \bsc{Rannou}, Clément \bsc{Guin}}
\date{17 mai 2016}
\begin{document}
\maketitle

%& \grave{a}
\section{Cahier des Charges}
\subsection{Presentation}
Le but de ce projet est de créer un clone du jeu vidéo \textit{Bomberman}. Il s'agit d'un jeu multijoueur qui se joue à 2 joueurs. Voici les régles :
\subsubsection{Regles du jeu}
Imaginez une arène (vue de dessus) composée de blocs incassables, de blocs cassables (briques) et de chemins praticables. Le jeu se joue à 2 joueurs, ce sont tous les 2 des Bombermans. Il peut se déplacer et poser des bombes. Après quelques secondes elles explosent et détruisent les blocs cassables proches. Les régles sont simples : il ne doit en rester qu'un. Il faut utiliser toute sa ruse, récuper des items bonus et en posant plusieurs bombes afin de venir à bout de votre adversaire.

\subsubsection{Environnement}
Nous décrivons dans cette section l'arène du jeu :
La construction de l'arène se fait à l'aide des images suivantes :
\begin{itemize}
\item Des blocs \textit{incassables} %image
\includegraphics[width=1cm,angle=0]{../blocs}
\item Des blocs \textit{cassables} %image
\includegraphics[width=1cm,angle=0]{../briques}
\item Des  2 personnages  %les 2 images
\includegraphics[width=1cm,angle=0]{../joueur1}
\includegraphics[width=1cm,angle=0]{../joueur2}
\end{itemize}
Il y aura aussi les items suivants dans l'arène :
\begin{itemize}
\item Les 2 bombes pour chaque personnage :
\includegraphics[width=1cm,angle=0]{../bomberouge}
\includegraphics[width=1cm,angle=0]{../bombebleue}
\item Le bonus de portée de la bombe
\includegraphics[width=1cm,angle=0]{../bonusbombe}
\item Le bonus de vie qui donne des points de vie 
\includegraphics[width=1cm,angle=0]{../vie}
\item Les 2 pièges pour chaque personnages
\includegraphics[width=1cm,angle=0]{../trap1}
\includegraphics[width=1cm,angle=0]{../trap2}
\end{itemize}

Chaque joueur commence à des coins opposés de l'arène. Ils peuvent se déplacer dans 4 directions, Nord, Sud, Est, Ouest. 	Ils ont chacun 3 vies en début de partie. Ils en perdent une s'ils sont dans la portée de la bombe. Si un joueur n'a plus de vie il a perdu. Chaque joueur peut poser 1 piège qui immobilise le joueur adverse
\subsubsection{Ce que le joueur peut ou doit faire}
\begin{itemize}
\item Ne rien faire
\item Se déplacer dans l'arène sur les cases vides
\item Pouvoir récuper des bonus
\item Aller sur la même case qu'un joueur, une bombe ou un bonus
\item Poser des bombes
\end{itemize}
\subsubsection{Ce que le joueur ne peut pas faire}
\begin{itemize}
\item Sortir de l'arène
\item Traverser des obstacles (blocs et briques)
\item Jouer une fois mort
\item Poser plus de 3 bombes en même temps
\end{itemize}
\subsubsection{Ce que les bombes peuvent ou doivent faire}
\begin{itemize}
\item Exploser au bout de 2 secondes
\item Être traversées par les joueurs
\end{itemize}
\subsubsection{Ce que les bombes ne peuvent pas faire}
\begin{itemize}
\item Cassser des blocs (éléments incassables)
\end{itemize}
\subsubsection{Ce que l'explosion peut ou doit faire}
\begin{itemize}
\item Tuer le joueur qui a posé la bombe ou/et le joueur adversaire
\item Casser des briques (éléments cassables) correspondant à la portée du joueur
\item Faire apparaître différents bonus
\end{itemize}
\subsubsection{Bonus}
\begin{itemize}
\item Vie
\item Augmentation de la portée de la bombe
\end{itemize}
\subsection{Fonctions}
Dans cette partie nous présenterons les différentes fonctions qui permettront le bon fonctionnement du jeu :
\begin{description}
\item [menu :] Fait apparaitre le menu
\item [startgame :] 	Démarre le jeu
\item [dessinermap :] Dessine l'arène de jeu
\item [ajouter\_brique :] Dessine une brique
\item [ajouter\_bloc :] Dessine un bloc
\item [bombe1 :] Pose une bombe et explose 2 secondes après pour le joueur 1
\item [bombe2 :] Pose une bombe et explose 2 secondes après pour le joueur 2
\item [explosion :] Fait exploser la bombe
\item [destruction\_animation\_explosion :] Détruit l'animation de la bombe
\item [bonus\_bombe :] Pose un bonus de portée de la bombe avec une probabilité de 5\%
\item [verif\_bonus\_bombe :] Verifie si le joueur 1 marche sur le bonus
\item [verif\_bonus\_bombe2 :] Verifie si le joueur 2 marche sur le bonus
\item [trap1 :] Pose le piège pour le joueur 1
\item [trap2 :] Pose le piège pour le joueur 2
\item [personnages :] Crée les 2 personnes
\item [animdroite :] Déplace le joueur 1 vers la droite
\item [animgauche :] Déplace le joueur 1 vers la gauche
\item [animbas :] Déplace le joueur 1 vers le bas
\item [animhaut :] Déplace le joueur 1 vers le haut
\item [animdroite2 :] Déplace le joueur 2 vers la droite
\item [animgauche2 :] Déplace le joueur 2 vers la gauche
\item [animbas2 :] Déplace le joueur 2 vers le bas
\item[animhaut2 :] Déplace le joueur 2 vers le haut
\end{description}
\section{Compte-Rendu}
\subsection{Recherche documentaire}
Passionnés d'informatique, nous connaissions le site openclassrooms (anciennement site du zéro), qui est un site de cours en ligne gratuit, orienté vers la programmation informatique. Nous avons donc suivi le cours d'initiation à python pour voir des notions que l'on avait pas vu en cours, ou réviser et approfondir celles qu'on avait vu en cours. 
\subsection{Travail d'équipe}
\subsubsection{Github}
Tout d'abord nous nous sommes heurté à des difficulter pour travailler ensemble, partager et mettre en communs nos travaux. Nous avons donc cherché une solution : Git. Git est un logiciel de gestion de versions, c'est à dire qu'il permet de stocker nos fichiers en conservants la chronologie de toutes les modifications qui ont été effectuées dessus. Pour une travail d'équipe efficace nous avons utilisé le service web git hub, qui est un réseau social utilisant le système Git. Git Hub nous a permis de s'organiser, ne pas perdre là trace de nos anciennes modifications. Par exemple, lorsque l'utilisateur fait une modification dans son code, et le code ne marche plus, on peut revenir à une version antérieur ou le code marchait pour pouvoir identifier le problème et le régler.
\subsection{Répartition du travail}
\subsubsection{Clément Guin}
\subsubsection{Titouan Rannou}
\subsubsection{Tanguy Thomas}
























			% {\Large } : pour écrire en plus gros 
			% \textit{} :  pour mettre en italiques
			% pour mettre en gras {\bf course}}
			% pour entourer un mot \fbox{la vie est belle}
\newpage 
\end{document}